\section{Система массового обслуживания M|M|c|K}

\begin{definition}{Система с ограниченной емкостью (M|M|c|K)}
Многоканальная система массового обслуживания, состоящая из $c$ идентичных приборов и буфера ограниченной емкости. Максимально допустимое число заявок в системе равно $K$ (из них $c$ могут находиться на обслуживании, а $K-c$ — в очереди). Если в момент поступления новой заявки в системе уже находится $N = K$ заявок, вновь прибывшая заявка получает отказ и покидает систему необслуженной.
\end{definition}

Интенсивность входного потока зависит от текущего состояния системы и задается следующим образом:
$$
\lambda_n = 
\begin{cases} 
\lambda, & 0 \le n < K, \\ 
0, & n \ge K.
\end{cases}
$$

Поскольку пространство состояний марковской цепи конечно ($N \in \{0, 1, \dots, K\}$), стационарное распределение вероятностей существует при любых значениях интенсивностей $\lambda$ и $\mu$. В отличие от систем с бесконечной очередью, ограничение $\rho < 1$ здесь не является обязательным: система сохраняет стабильность даже при сверхвысоких нагрузках за счет отбрасывания избыточных заявок.

\begin{center}
    \begin{tikzpicture}[
        scale=1.5,
        >={Stealth[length=3mm]},
        state/.style={circle, draw, minimum size=0.9cm, thick},
        every node/.style={font=\small}
    ]
        % Состояния
        \node[state] (skm1) at (-2,0) {$K-1$};
        \node[state] (sk) at (0,0) {$K$};
        \node[state, dashed, text=gray, draw=gray] (sk1) at (2.5,0) {$K+1$};

        % Переходы до K
        \draw[->, thick] (-3.5, 0.2) to[bend left=20] node[midway, above] {$\lambda$} (skm1);
        \draw[<-, thick] (-3.5, -0.2) to[bend right=20] node[midway, below] {$c\mu$} (skm1);

        \draw[->, thick] (skm1) to[bend left=30] node[midway, above] {$\lambda$} (sk);
        \draw[<-, thick] (skm1) to [bend right=30] node[midway, below] {$c\mu$} (sk);

        % Запрещенные переходы за K
        \path[draw, ->, thick, gray, dashed] (sk) to [bend left=30] coordinate[midway] (cross1) (sk1);
        \path[draw, <-, thick, gray, dashed] (sk) to[bend right=30] coordinate[midway] (cross2) (sk1);
        
        % Красные кресты на переходах
        \draw[red, ultra thick] ($(cross1)+(-0.15, -0.15)$) -- ($(cross1)+(0.15, 0.15)$);
        \draw[red, ultra thick] ($(cross1)+(-0.15, 0.15)$) -- ($(cross1)+(0.15, -0.15)$);
        
        \draw[red, ultra thick] ($(cross2)+(-0.15, -0.15)$) -- ($(cross2)+(0.15, 0.15)$);
        \draw[red, ultra thick] ($(cross2)+(-0.15, 0.15)$) -- ($(cross2)+(0.15, -0.15)$);

        \node[red, font=\footnotesize] at (1.25, 0.8) {$\lambda_K = 0$};
    \end{tikzpicture}
    \captionof*{figure}{\small Фрагмент графа состояний системы M|M|c|K. Переходы в состояния $n > K$ заблокированы.}
\end{center}

\begin{statement}[Стационарное распределение системы M|M|c|K]
Для многоканальной системы с ограничением емкости $K$ вероятности состояний задаются формулами:
$$
p_n = 
\begin{cases} 
\displaystyle \frac{r^n}{n!} p_0, & 0 \le n < c, \\
\displaystyle \frac{r^n}{c^{n-c} c!} p_0, & c \le n \le K,
\end{cases}
$$
где $r = \frac{\lambda}{\mu}$ — приведенная нагрузка, $\rho = \frac{r}{c}$ — коэффициент загрузки. Вероятность простоя системы вычисляется как:
$$
p_0 = \left( \sum_{n=0}^{c-1} \frac{r^n}{n!} + \frac{r^c}{c!} S_K(\rho) \right)^{-1}, \quad \text{где} \quad S_K(\rho) = 
\begin{cases} 
\displaystyle \frac{1 - \rho^{K-c+1}}{1 - \rho}, & \rho \ne 1, \\ 
K - c + 1, & \rho = 1.
\end{cases}
$$
\end{statement}

\begin{sproof}
Решение уравнений глобального баланса для процессов рождения и гибели дает стандартное выражение вероятностей через $p_0$. Нормировочная константа $p_0$ находится из условия $\sum_{n=0}^K p_n = 1$. Сумма геометрической прогрессии $S_K(\rho) = \sum_{n=c}^{K} \rho^{n-c}$ вычисляется по стандартной формуле суммы конечной геометрической прогрессии при $\rho \ne 1$, а при $\rho = 1$ сводится к сумме $(K - c + 1)$ единиц.
\end{sproof}

\begin{note}
Ключевым отличием систем с отказами от классических систем является понятие \textbf{эффективной интенсивности входного потока} $\lambda_{\text{eff}}$. Поскольку заявки, поступающие в момент, когда система находится в состоянии $K$ (с вероятностью $p_K$), отбрасываются, фактический поток заявок, принимаемых к обслуживанию, меньше номинального $\lambda$. Следовательно:
$$
\lambda_{\text{eff}} = \lambda (1 - p_K).
$$
Все глобальные макроскопические характеристики системы (в частности, закон Литтла) обязаны строиться исключительно с использованием $\lambda_{\text{eff}}$, так как отброшенные заявки не задерживаются в системе.
\end{note}

\begin{statement}[Макроскопические характеристики системы M|M|c|K]
Для макроскопических метрик эффективности справедливы следующие соотношения:
\begin{itemize}
    \item Среднее число заявок в очереди: 
    $$ L_q = \begin{cases}
    \displaystyle \frac{p_0 r^c \rho}{c! (1 - \rho)^2} \Big[ 1 - \rho^{K-c+1} - (1 - \rho) (K - c + 1) \rho^{K-c} \Big], & \rho \ne 1, \\
    \displaystyle \frac{p_0 r^c}{c!} \cdot \frac{(K - c)(K - c + 1)}{2}, & \rho = 1.
    \end{cases} $$
    \item Среднее число заявок в системе: $\displaystyle L = L_q + \frac{\lambda_{\text{eff}}}{\mu}$.
    \item Среднее время пребывания в системе: $\displaystyle W = \frac{L}{\lambda_{\text{eff}}} = \frac{L}{\lambda(1 - p_K)}$.
    \item Среднее время ожидания в очереди: $\displaystyle W_q = \frac{L_q}{\lambda_{\text{eff}}} = W - \frac{1}{\mu}$.
\end{itemize}
\end{statement}

\begin{sproof}
\textbf{1. Среднее число заявок в очереди ($L_q$).} По определению математического ожидания имеем:
$$
L_q = \sum_{n=c+1}^{K} (n - c) p_n = \frac{p_0 r^c}{c!} \sum_{n=c+1}^{K} (n - c) \rho^{n-c}.
$$
Введем новый индекс суммирования $j = n - c$. Тогда граница суммирования изменится от $j = 1$ до $N = K - c$:
$$
L_q = \frac{p_0 r^c}{c!} \sum_{j=1}^{N} j \rho^j.
$$
Рассмотрим случай $\rho \ne 1$. Применим метод дифференцирования по параметру $\rho$:
$$
\sum_{j=1}^{N} j \rho^j = \rho \sum_{j=1}^{N} j \rho^{j-1} = \rho \frac{d}{d\rho} \left( \sum_{j=1}^{N} \rho^j \right) = \rho \frac{d}{d\rho} \left( \frac{\rho - \rho^{N+1}}{1 - \rho} \right).
$$
Вычислим производную дроби:
$$
\frac{d}{d\rho} \left( \frac{\rho - \rho^{N+1}}{1 - \rho} \right) = \frac{(1 - (N+1)\rho^N)(1 - \rho) - (\rho - \rho^{N+1})(-1)}{(1 - \rho)^2} = \frac{1 - (N+1)\rho^N + N\rho^{N+1}}{(1 - \rho)^2}.
$$
Умножая на $\rho$ и возвращаясь к исходным обозначениям ($N = K - c$), преобразуем числитель:
$$
1 - (K - c + 1)\rho^{K-c} + (K - c)\rho^{K-c+1} = 1 - \rho^{K-c+1} - (1 - \rho)(K - c + 1)\rho^{K-c}.
$$
Подставляя это в выражение для $L_q$, получаем требуемую формулу для $\rho \ne 1$.

Для случая $\rho = 1$ ряд представляет собой сумму арифметической прогрессии:
$$
\sum_{j=1}^{K-c} j \cdot 1^j = \frac{(K - c)(K - c + 1)}{2}, \quad \text{откуда} \quad L_q = \frac{p_0 r^c}{c!} \frac{(K - c)(K - c + 1)}{2}.
$$

\textbf{2. Временные характеристики.} Формулы для $L, W$ и $W_q$ прямо следуют из применения универсального закона Литтла к системе с отбрасыванием заявок. В качестве интенсивности потока выступает величина $\lambda_{\text{eff}}$.
\end{sproof}

\begin{conseq}
В частном случае одноканальной системы с ограничением ($c=1$, система M|M|1|K) уравнение глобального баланса интенсивностей принимает наглядный вид. Среднее число заявок $L = L_q + (1 - p_0)$, где $(1 - p_0)$ — вероятность того, что прибор занят. Из соотношения $L = L_q + \lambda_{\text{eff}}/\mu$ немедленно следует:
$$
\underbrace{\mu (1 - p_0)}_{\substack{\text{эффективная интенсивность} \\ \text{обслуживания}}} = \underbrace{\lambda (1 - p_K)}_{\substack{\text{эффективная интенсивность} \\ \text{входящего потока}}}.
$$
Система достигает равновесия, при котором средняя скорость обработки находящихся в приборе заявок строго равна скорости поступления новых (принятых) заявок.
\end{conseq}

\begin{example}{Аналитический расчет вероятности отказа для системы M|M|1|K}
Рассмотрим одноканальную систему с ограниченным буфером (общая емкость системы равна $K$, число приборов $c=1$). Продемонстрируем вывод явной аналитической зависимости вероятности потери заявки $p_K$ от коэффициента загрузки $\rho$ и исследуем ее поведение при высоких нагрузках ($\rho \ge 1$).
\end{example}

\begin{eproof}
Для одноканальной системы $c=1$, следовательно, множители $c! = 1$ и $c^{n-c} = 1$. Общая формула для стационарных вероятностей $p_n$ существенно упрощается на всем диапазоне состояний $0 \le n \le K$:
$$
p_n = \rho^n p_0.
$$
Используя условие нормировки $\sum_{n=0}^K p_n = 1$, находим константу $p_0$ через сумму конечной геометрической прогрессии. Для $\rho \ne 1$ получаем:
$$
p_0 \sum_{n=0}^K \rho^n = p_0 \frac{1 - \rho^{K+1}}{1 - \rho} = 1 \quad \implies \quad p_0 = \frac{1 - \rho}{1 - \rho^{K+1}}.
$$
Вероятность того, что система полностью заполнена (и поступающая заявка получит отказ), вычисляется подстановкой $n = K$:
$$
p_K = \rho^K p_0 = \frac{\rho^K (1 - \rho)}{1 - \rho^{K+1}} = \frac{\rho^K (\rho - 1)}{\rho^{K+1} - 1}.
$$
В отличие от систем с бесконечной очередью, где при $\rho \to 1$ длина очереди стремится к бесконечности (асимптота), система M|M|1|K сохраняет стабильность при любых нагрузках. Рассмотрим два критических режима:

\begin{itemize}
    \item \textbf{Пограничная загрузка ($\rho = 1$):} Возникает неопределенность вида $\left[ \frac{0}{0} \right]$. Раскрывая ее по правилу Лопиталя (или подставляя $\rho=1$ напрямую в исходную сумму $\sum_{n=0}^K 1^n = K+1$), получаем строго конечную вероятность отказа:
    $$
    p_K \Big|_{\rho=1} = \frac{1}{K+1}.
    $$
    
    \item \textbf{Режим перегрузки ($\rho > 1$):} При интенсивностях входного потока, превышающих скорость обслуживания, очередь не возрастает бесконечно, а стабилизируется за счет отбрасывания лишних заявок. Разделив числитель и знаменатель дроби на $\rho^K$, получим:
    $$
    p_K = \frac{\rho - 1}{\rho - \rho^{-K}}.
    $$
    При сильной перегрузке ($\rho \gg 1$) экспоненциально убывающим слагаемым $\rho^{-K}$ можно пренебречь. Это дает следующую асимптотику вероятности отказа:
    $$
    p_K \approx \frac{\rho - 1}{\rho} = 1 - \frac{1}{\rho}.
    $$
\end{itemize}

Данный математический результат имеет строгий физический смысл. Вычислим эффективную интенсивность потока, проходящего через систему в режиме перегрузки:
$$
\lambda_{\text{eff}} = \lambda (1 - p_K) \approx \lambda \left( 1 - \left( 1 - \frac{1}{\rho} \right) \right) = \lambda \cdot \frac{1}{\rho} = \lambda \cdot \frac{\mu}{\lambda} = \mu.
$$
Система автоматически адаптируется к любой перегрузке: она работает на абсолютном пределе своей производительности (обслуживая $\mu$ заявок в единицу времени), а весь избыточный трафик отбрасывается на входе.
\end{eproof}

\begin{center}
\begin{tikzpicture}
    \begin{axis}[
        width=0.8\textwidth, height=7.5cm,
        xlabel={Коэффициент загрузки $\rho$}, 
        ylabel={Вероятность отказа $p_K$},
        xmin=0, xmax=2.5, ymin=0, ymax=0.8,
        grid=major,
        legend pos=north west,
        samples=100, thick,
        cycle list={
            {Red3, solid},
            {DeepSkyBlue4, dashed},
            {SpringGreen4, dotted},
            {black, dashdotted}
        }
    ]
        % Формула для K=2
        \addplot+[domain=0.01:2.5] {x^2*(1-x)/(1-x^3)};
        \addlegendentry{$K=2$}
        
        % Формула для K=5
        \addplot+[domain=0.01:2.5] {x^5*(1-x)/(1-x^6)};
        \addlegendentry{$K=5$}
        
        % Формула для K=10
        \addplot+[domain=0.01:2.5] {x^10*(1-x)/(1-x^11)};
        \addlegendentry{$K=10$}
        
        % Асимптота 1 - 1/\rho для \rho > 1
        \addplot+[domain=1.001:2.5] {1 - 1/x};
        \addlegendentry{Асимптота $1 - \rho^{-1}$}
        
    \end{axis}
\end{tikzpicture}
\captionof*{figure}{\small Зависимость вероятности потери заявки $p_K$ от загрузки $\rho$ для различной емкости буфера. Черной штрихпунктирной линией показана предельная асимптота при $K \to \infty$ в режиме перегрузки ($\rho > 1$).}
\end{center}

\subsection{Распределение состояний в моменты поступления заявок}

В системах с ограничением емкости буфера необходимо строго различать три типа распределений:
\begin{enumerate}
    \item $p_n$ — стационарное распределение (доля времени, которую система проводит в состоянии $n$).
    \item $q_n$ — распределение вероятностей состояний, которые застает \textit{произвольная} поступающая заявка.
    \item $q_n^*$ — распределение вероятностей состояний, которые застает \textit{успешно принятая} в систему заявка.
\end{enumerate}

Согласно свойству PASTA (Poisson Arrivals See Time Averages), пуассоновский поток заявок «видит» систему в стационарном состоянии, поэтому $q_n = p_n$. Однако для систем $M|M|c|K$ факт приема заявки зависит от ее состояния: если $n = K$, заявка будет отброшена. Следовательно, распределение успешно принятых заявок $q_n^*$ является \textbf{условным} распределением и не совпадает с $q_n$.

\begin{note}
Напомним \textbf{формулу Байеса} для условной вероятности гипотез. Если $H_i$ задает полную группу несовместных гипотез, а $A$ является некоторым событием, то апостериорная вероятность гипотезы $H_n$ вычисляется как:
$$
\P(H_n \mid A) = \frac{\P(A \mid H_n) \P(H_n)}{\dsum_{i} \P(A \mid H_i) \P(H_i)}.
$$
\end{note}

\begin{statement}[Связь вероятностей $p_n$ и  $q_n^*$]
Вероятность того, что успешно принятая в систему $M|M|c|K$ заявка застанет в ней ровно $n$ заявок, определяется выражением:
$$
q_n^* = 
\begin{cases}
\displaystyle \frac{p_n}{1 - p_K}, & 0 \le n \le K-1, \\
0, & n = K.
\end{cases}
$$
\end{statement}

\begin{sproof}
Зафиксируем произвольный момент времени $t$ в стационарном режиме. Рассмотрим малый интервал $(t, t + \Delta t)$. 
Пусть гипотеза $H_n$ заключается в том, что в момент $t$ система находится в состоянии $n$. Вероятности этих гипотез по определению суть стационарные вероятности: $\P(H_n) = p_n$.

Определим событие $A$ как «на интервале $(t, t + \Delta t)$ в систему пришла и была \textbf{успешно принята} новая заявка». Согласно определению условной вероятности:
$$
\P(A \mid H_n) = \frac{\P(A \cap H_n)}{\P(H_n)}.
$$

Рассмотрим числитель для случаев $n \in \{0, \dots, K-1\}$. В этих состояниях система имеет хотя бы одно свободное место, поэтому любая пришедшая заявка принимается с вероятностью $1$. Следовательно, при $n < K$ пересечение $A \cap H_n$ совпадает с пересечением $\{\text{на интервале } \Delta t \text{ пришла 1 заявка}\} \cap H_n$. Используя свойства пуассоновского потока, получаем:
$$
\P(A \mid H_n) = \frac{\P(\{\text{пришла 1 заявка}\} \cap H_n)}{\P(H_n)} = \frac{(\lambda \Delta t + o(\Delta t)) \cdot p_n}{p_n} = \lambda \Delta t + o(\Delta t).
$$

Для состояния $n = K$ система полностью заполнена. Любая пришедшая заявка получает отказ, то есть событие $A$ в состоянии $K$ невозможно: $A \cap H_K = \varnothing$, откуда $\P(A \mid H_K) = 0$.

Запишем условные вероятности в виде системы:
$$
\P(A \mid H_n) = 
\begin{cases}
\lambda \Delta t + o(\Delta t), & 0 \le n \le K-1, \\
0, & n = K.
\end{cases}
$$

Воспользуемся формулой Байеса для нахождения искомой вероятности $q_n^*$:
$$
q_n^* \doteq \lim_{\Delta t \to 0} \P(H_n \mid A) = \lim_{\Delta t \to 0} \frac{\P(A \mid H_n) \P(H_n)}{\dsum_{i=0}^{K} \P(A \mid H_i) \P(H_i)} = \lim_{\Delta t \to 0} \frac{\P(A \mid H_n) \cdot p_n}{\dsum_{i=0}^{K-1} (\lambda \Delta t + o(\Delta t)) p_i + 0 \cdot p_K}.
$$

Для любого $n \le K-1$ подставим значение $\P(A \mid H_n)$:
$$
q_n^* = \lim_{\Delta t \to 0} \frac{(\lambda \Delta t + o(\Delta t)) p_n}{(\lambda \Delta t + o(\Delta t)) \dsum_{i=0}^{K-1} p_i} = \lim_{\Delta t \to 0} \frac{\left(\lambda + \frac{o(\Delta t)}{\Delta t}\right) p_n}{\left(\lambda + \frac{o(\Delta t)}{\Delta t}\right) \dsum_{i=0}^{K-1} p_i}.
$$

Переходя к пределу при $\Delta t \to 0$, находим:
$$
q_n^* = \frac{\lambda p_n}{\lambda \dsum_{i=0}^{K-1} p_i} = \frac{p_n}{\dsum_{i=0}^{K-1} p_i}.
$$
В силу нормировочного условия $\dsum_{i=0}^{K} p_i = 1$, сумма в знаменателе равна $1 - p_K$. Следовательно:
$$
q_n^* = \frac{p_n}{1 - p_K}, \quad n \in \{0, \dots, K-1\}.
$$
Для состояния $n = K$ числитель под знаком предела тождественно равен нулю в силу $\P(A \mid H_K) = 0$, откуда $q_K^* = 0$.
\end{sproof}

\begin{conseq}
Из полученной формулы следует соотношение $q_n^* (1 - p_K) = p_n$. При устремлении емкости системы к бесконечности ($K \to \infty$) мы переходим к классической системе $M|M|c|\infty$. В этом случае вероятность отказа стремится к нулю ($p_K \to 0$), и знаменатель дроби стремится к единице. Следовательно, $q_n^* \to p_n$. Аналитически это означает, что в системах с бесконечной очередью все заявки принимаются, и условное распределение совпадает с безусловным стационарным распределением.
\end{conseq}

\begin{note}
Поскольку для процессов с непрерывным временем вероятность поступления заявки в строго зафиксированный момент $t$ равна нулю, прямое вычисление условной вероятности \linebreak $\P(H_n \mid \{ \text{поступление заявки в момент } t \})$ невозможно. Введение предельного перехода при $\Delta t \to 0$ позволяет корректно определить условную вероятность относительно события нулевой меры и формализовать понятие состояния системы непосредственно перед моментом прибытия заявки.
\end{note}

\subsection{Системы с отказами (M|M|c|c) и формула Эрланга}

Особый теоретический и практический интерес (в частности, в телефонии и инфокоммуникационных сетях) представляет вырожденный случай системы с ограниченной емкостью, в которой места для ожидания полностью отсутствуют.

\begin{definition}{Система с явными потерями (M|M|c|c)}
Многоканальная система массового обслуживания, состоящая из $c$ приборов, в которой емкость буфера равна нулю. Максимально допустимое число заявок в системе в точности совпадает с числом приборов ($K = c$). Любая заявка, заставшая все $c$ каналов занятыми, мгновенно получает отказ и покидает систему.
\end{definition}

\begin{statement}[Первая формула Эрланга]
Для системы $M|M|c|c$ стационарное распределение вероятностей состояний задается формулой:
$$
p_n = \frac{\displaystyle \frac{r^n}{n!}}{\displaystyle \sum_{i=0}^c \frac{r^i}{i!}}, \quad n \in \{0, 1, \dots, c\},
$$
где $r = \frac{\lambda}{\mu}$ — приведенная нагрузка. 
Вероятность того, что все каналы заняты (вероятность блокировки или потери заявки), обозначается как $B(c, r)$ и называется \textbf{формулой потерь Эрланга} (Erlang-B formula):
$$
B(c, r) \equiv p_c = \frac{\displaystyle \frac{r^c}{c!}}{\displaystyle \sum_{i=0}^c \frac{r^i}{i!}}.
$$
\end{statement}

\begin{sproof}
Поскольку $K = c$, состояния с наличием очереди ($n > c$) не существуют. Воспользуемся общим решением для системы $M|M|c|K$. При $n \le c$ интенсивность обслуживания линейна и равна $n\mu$. Уравнения глобального баланса дают $p_n = \frac{r^n}{n!} p_0$. 

Применяя условие нормировки $\dsum_{n=0}^c p_n = 1$, вычисляем вероятность простоя системы:
$$
p_0 = \left[ \dsum_{i=0}^c \frac{r^i}{i!} \right]^{-1}.
$$
Подстановка $p_0$ в выражение для $p_n$ напрямую приводит к искомой формуле. Полагая $n = c$, получаем выражение для $B(c, r)$.
\end{sproof}

\begin{note}[1]
    Стационарное распределение $p_n$ обладает свойством инвариантности по отношению к форме распределения времени обслуживания. Согласно теореме Севастьянова, формулы Эрланга остаются справедливыми для системы общего вида $M|G|c|c$. Вероятности состояний зависят от распределения времени обслуживания исключительно через его математическое ожидание $\E[S]$ (посредством параметра нагрузки $r = \lambda \E[S]$).
\end{note}    

\begin{note}[2]
    Вероятность блокировки заявки тождественно совпадает со стационарной вероятностью $p_c$ в силу свойства PASTA. Данное свойство применимо к исходному пуассоновскому потоку всех поступающих заявок (как успешно принимаемых системой, так и получающих отказ), в отличие от условного потока только принятых заявок.
\end{note}

\begin{note}[3]
    Прямое вычисление вероятности отказа по формуле Эрланга $B(c, r) = p_c$ при больших значениях $c$ приводит к вычислительным трудностям из-за наличия факториалов. Для численного нахождения стационарных вероятностей применяют рекуррентное соотношение:
    $$
    B(c, r) = \frac{r B(c-1, r)}{c + r B(c-1, r)}, \quad B(0, r) = 1.
    $$
\end{note}

\section{Система массового обслуживания M|G|1 и формула Полячека — Хинчина}

Рассмотрим одноканальную систему массового обслуживания с пуассоновским входным потоком заявок и произвольным (General) распределением времени обслуживания. Дисциплина обслуживания — FCFS (First Come, First Served). 

\begin{definition}{Система с произвольным обслуживанием (M|G|1)}
Система, в которой интенсивность входного пуассоновского потока равна $\lambda$, а времена обслуживания отдельных заявок $\{S_1, S_2, \dots, S_n, \dots\}$ образуют последовательность независимых одинаково распределенных случайных величин с функцией распределения $B(t) = \P(S_n \le t)$. Среднее время обслуживания и коэффициент загрузки системы определяются как:
$$
\mu \doteq \frac{1}{\E[S]}, \quad \rho \doteq \frac{\lambda}{\mu} < 1.
$$
\end{definition}

\begin{note}[1]
Поскольку время обслуживания может иметь как непрерывный, так и дискретный характер (или их комбинацию), математическое ожидание $\E[S]$ строго определяется через \textbf{интеграл Лебега — Стилтьеса}:
$$
\E[S] = \int_{0}^{+\infty} t \, dB(t).
$$
Данный интеграл обобщает классический интеграл Римана и сумму для дискретных величин, выполняя интегрирование по вероятностной мере, порожденной функцией распределения $B(t)$. В случае абсолютно непрерывного распределения с плотностью $b(t)$ он сводится к $\int t b(t) dt$.
\end{note}

\begin{note}[2]
Требование независимости и одинаковой распределенности величин $S_n$ является критически важным. Если в реальной системе присутствует физический износ оборудования (время обработки растет с номером заявки $n$) или адаптивность (скорость обслуживания зависит от текущей длины очереди $X_n$), распределение $S_{n+1}$ начинает зависеть от предыстории, делая приведенные ниже выкладки непригодными.
\end{note}

Для анализа классической системы $M|G|1$ применение стандартного аппарата уравнений Колмогорова напрямую невозможно. Случайный процесс $X(t)$, описывающий число заявок в системе в момент времени $t$, не обладает марковским свойством: вероятность перехода в новое состояние зависит не только от текущего числа заявок, но и от того, \textit{сколько времени} уже обслуживается текущая заявка (поскольку произвольное распределение времени обслуживания не обладает свойством отсутствия памяти, в отличие от экспоненциального). 

Для преодоления этой проблемы Дэвид Кендалл предложил метод дискретизации времени, при котором эволюция системы рассматривается исключительно в строго определенные, «регенеративные» моменты.

\begin{definition}{Вложенная цепь Маркова (метод Кендалла)}
Пусть $t_n$ — моменты окончания обслуживания $n$-й заявки. Рассмотрим дискретный процесс $X_n \doteq X(t_n+0)$, описывающий число заявок в системе \textit{непосредственно после} ухода $n$-й заявки. Последовательность случайных величин $\{X_n\}_{n=1}^{\infty}$ образует дискретную однородную цепь Маркова, называемую \textbf{вложенной}.
\end{definition}

\begin{note}[1]
В классических марковских цепях с дискретным временем индекс $n$ обычно соответствует детерминированным шагам (например, фиксированным интервалам $\Delta t$). В методе Кендалла физическое время $t_n(\omega)$ окончания обслуживания $n$-й заявки является случайной величиной, зависящей от конкретной реализации $\omega$. 

Корректность построения цепи обеспечивается сменой индексного множества. Параметр $n \in \mathbb{N}$ здесь выступает не как физическое время, а как \textbf{порядковый номер ухода $n$-й заявки}. Переход от непрерывного процесса $X(t)$ к дискретной последовательности $\{X_n\}_{n=1}^{\infty}$ абстрагируется от того, \textit{когда именно} произошли переходы. Поскольку индексы $n$ пробегают фиксированное детерминированное множество $\mathbb{N}$, $\{X_n\}$ представляет собой корректно определенный случайный процесс с дискретным временем.
\end{note}

\begin{note}[2]
Строгое доказательство марковости вложенной цепи опирается на алгебраическое уравнение эволюции:
$$
X_{n+1} = \max(0, X_n - 1) + A_{n+1},
$$
где $A_{n+1}$ — число заявок, поступивших в систему за время $S_{n+1}$ (обслуживание $(n+1)$-й заявки). 

\begin{center}
    \begin{tikzpicture}[
        scale=1.2,
        >={Stealth[length=2.5mm]},
        every node/.style={font=\small},
        % Подключаем библиотеку декораций в преамбуле: \usetikzlibrary{decorations.pathreplacing, arrows.meta}
    ]
        % Ось времени
        \draw[->, thick] (0,0) -- (10,0) node[right] {$t$};
        
        % Моменты прихода (стрелки вниз)
        \draw[<-, DeepSkyBlue4, thick] (1.5,0.1) -- (1.5,1) node[above] {Новые заявки};
        \draw[<-, DeepSkyBlue4, thick] (3.8,0.1) -- (3.8,1);
        \draw[<-, DeepSkyBlue4, thick] (5.2,0.1) -- (5.2,1);
        \draw[<-, DeepSkyBlue4, thick] (7.9,0.1) -- (7.9,1);
        
        % Моменты ухода (стрелки вверх и узлы t_n)
        % Смещаем t_{n-1} чуть левее, чтобы не мешать скобке
        \draw[->, Red3, thick] (2.5,-1.2) node[left] {$t_{n-1}$} -- (2.5,-0.1);
        \filldraw (2.5,0) circle (1.5pt);
        \node[above right, xshift=2pt] at (2.5, 0) {$X_{n-1}$};

        \draw[->, Red3, thick] (6.5,-1.2) node[left] {$t_n$} -- (6.5,-0.1);
        \filldraw (6.5,0) circle (1.5pt);
        \node[above right, xshift=2pt] at (6.5, 0) {$X_n$};
        
        \draw[->, Red3, thick] (9.0,-1.2) node[right] {$t_{n+1}$} -- (9.0,-0.1);
        \filldraw (9.0,0) circle (1.5pt);
        \node[above right, xshift=2pt] at (9.0, 0) {$X_{n+1}$};

        % Интервалы обслуживания (опущены ниже и увеличен отступ текста)
        \draw[decorate, decoration={brace, amplitude=5pt, mirror}] (2.6,-1.4) -- (6.4,-1.4) 
            node[midway, below=8pt] {$S_n$};
        
        \draw[decorate, decoration={brace, amplitude=5pt, mirror}] (6.6,-1.4) -- (8.9,-1.4) 
            node[midway, below=8pt] {$S_{n+1}$};
        
        % Поступившие заявки A_n
        \draw[decorate, decoration={brace, amplitude=5pt}] (6.5, 1.2) -- (9.0, 1.2) 
            node[midway, above=8pt] {$A_{n+1}$ заявок прибыло};
            
    \end{tikzpicture}
    \captionof*{figure}{\small Временная диаграмма процесса обслуживания. Марковская цепь формируется путем фиксации состояний $X_n$ строго в моменты ухода заявок.}
\end{center}

Поскольку моменты наблюдения $t_n$ строго привязаны к окончанию обслуживания, в момент начала обработки следующей заявки ее «таймер» стартует с нуля. Количество прибывших заявок $A_{n+1}$ зависит исключительно от длительности интервала $S_{n+1}$ и интенсивности пуассоновского потока $\lambda$. В силу фундаментальной аксиомы о том, что времена $\{S_n\}$ независимы от предыстории системы, величина $A_{n+1}$ также стохастически независима от прошлых состояний $X_{n-1}, X_{n-2}, \dots$.

Таким образом, распределение случайной величины $X_{n+1}$ полностью и однозначно определяется текущим состоянием $X_n$:
$$
\P(X_{n+1} = j \mid X_n, X_{n-1}, \dots, X_0) = \P(X_{n+1} = j \mid X_n).
$$
Однородность данной цепи (независимость переходных вероятностей от шага $n$) вытекает из стационарности пуассоновского потока и идентичности распределений $B(t)$ для всех заявок.
\end{note}

\begin{theorem}{Формула Полячека — Хинчина}
Для системы $M|G|1$ стационарное распределение вероятностей состояний существует тогда и только тогда, когда коэффициент загрузки $\rho < 1$. В стационарном режиме среднее число заявок в системе определяется формулой:
$$
L = \rho + \frac{\rho^2 + \lambda^2 \sigma_B^2}{2(1-\rho)},
$$
где $\sigma_B^2 = \mathbb{D}[S]$ — дисперсия времени обслуживания.
\end{theorem}

\begin{tproof}
Пусть $A_{n+1}$ — число заявок, поступивших в систему за время обслуживания $(n+1)$-й заявки. Эволюция состояний вложенной цепи Маркова описывается разностным стохастическим уравнением:
$$
X_{n+1} = 
\begin{cases} 
X_n - 1 + A_{n+1}, & X_n \ge 1, \\ 
A_{n+1}, & X_n = 0.
\end{cases}
$$
Введем индикаторную функцию события $\{X_n > 0\}$, обозначаемую $\mathbb{I}_{\{X_n > 0\}}$, равную $1$ при $X_n > 0$ и $0$ при $X_n = 0$. Уравнение эволюции принимает единый алгебраический вид:
\begin{equation}\label{eq:main}
    X_{n+1} = X_n - \mathbb{I}_{\{X_n > 0\}} + A_{n+1} \tag{$*$}
\end{equation}

Предположим, что марковская цепь эргодична. В стационарном режиме распределения $X_{n+1}$ и $X_n$ совпадают, следовательно, их начальные моменты равны. Обозначим искомое математическое ожидание числа заявок в момент ухода как $\E[X_{n+1}] = \E[X_n] \doteq L^{(D)}$. Верхний индекс $(D)$ (от англ. \textit{Departure} — уход) подчеркивает, что данная величина есть среднее число заявок в системе, наблюдаемое исключительно в моменты завершения обслуживания. В общем случае она может не совпадать со средним числом заявок в произвольный момент времени $L$.

\textbf{1. Первый начальный момент и вероятность непустой системы.}
Применим оператор математического ожидания к обеим частям уравнения \eqref{eq:main}:
$$
L^{(D)} = L^{(D)} - \E[\mathbb{I}_{\{X_n > 0\}}] + \E[A_{n+1}] \quad \implies \quad \E[\mathbb{I}_{\{X_n > 0\}}] = \E[A_{n+1}].
$$
По свойствам индикаторной функции ее математическое ожидание равно вероятности соответствующего события: $\E[\mathbb{I}_{\{X_n > 0\}}] = \P(X_n > 0) = 1 - \pi_0$, где $\pi_0$ — стационарная вероятность того, что система пуста в моменты ухода. Иными словами, это вероятность того, что уходящая заявка оставляет после себя непустую систему.

Вычислим $\E[A_{n+1}]$, применяя формулу полного математического ожидания по длительности обслуживания $S_{n+1}$. В силу свойств пуассоновского потока условное распределение $\{A_{n+1} \mid S_{n+1} = t\} \sim \text{Poiss}(\lambda t)$:
$$
\E[A_{n+1}] = \int_{0}^{+\infty} \E[A_{n+1} \mid S_{n+1} = t] \, dB(t) = \int_{0}^{+\infty} \lambda t \, dB(t) = \lambda \E[S_{n+1}] = \rho.
$$
Отсюда немедленно следует важный результат: $1 - \pi_0 = \rho$. Поскольку для существования стационарного режима вероятность пустого состояния должна быть строго положительной ($\pi_0 > 0$), из этого тождества аналитически вытекает \textbf{необходимое условие стационарности: $\rho < 1$}. (В пункте 4 мы строго докажем равенство распределений по временам ухода и в произвольный момент времени $\pi_0 = p_0$. Это придаст величине $\rho$ классический физический смысл — доли времени занятости прибора в непрерывном времени).

\textbf{2. Второй начальный момент и длина очереди.}
Возведем уравнение \eqref{eq:main} в квадрат. Учитывая идемпотентность индикатора $\mathbb{I}_{\{X_n > 0\}}^2 = \mathbb{I}_{\{X_n > 0\}}$ и свойство $X_n \mathbb{I}_{\{X_n > 0\}} = X_n$, получаем:
\begin{align*}
X_{n+1}^2 &= \big(X_n - \mathbb{I}_{\{X_n > 0\}} + A_{n+1}\big)^2 \\
&= X_n^2 + \mathbb{I}_{\{X_n > 0\}}^2 + A_{n+1}^2 - 2X_n\mathbb{I}_{\{X_n > 0\}} - 2A_{n+1}\mathbb{I}_{\{X_n > 0\}} + 2A_{n+1}X_n \\
&= X_n^2 + \mathbb{I}_{\{X_n > 0\}} + A_{n+1}^2 - 2X_n - 2A_{n+1}\mathbb{I}_{\{X_n > 0\}} + 2A_{n+1}X_n.
\end{align*}
Перейдем к математическим ожиданиям. В силу стационарности $\E[X_{n+1}^2] = \E[X_n^2]$. Приращение пуассоновского потока $A_{n+1}$ зависит только от длительности $S_{n+1}$ и строго не зависит от предыстории процесса (состояния $X_n$ и индикатора $\mathbb{I}_{\{X_n > 0\}}$). Следовательно, математическое ожидание произведений распадается на произведение математических ожиданий:
\begin{align*}
0 &= \E[\mathbb{I}_{\{X_n > 0\}}] + \E[A_{n+1}^2] - 2\E[X_n] - 2\E[A_{n+1}]\E[\mathbb{I}_{\{X_n > 0\}}] + 2\E[A_{n+1}]\E[X_n].
\end{align*}
Подставляя найденные ранее значения ($\E[\mathbb{I}_{\{X_n > 0\}}] = \rho$, $\E[A_{n+1}] = \rho$, $\E[X_n] = L^{(D)}$), приводим уравнение к виду:
$$
0 = \rho + \E[A_{n+1}^2] - 2L^{(D)} - 2\rho^2 + 2\rho L^{(D)}.
$$
Группируя слагаемые с $L^{(D)}$, находим:
$$
2L^{(D)}(1 - \rho) = \rho - 2\rho^2 + \E[A_{n+1}^2].
$$
Поскольку доказанное ранее условие стационарности $\rho < 1$ выполнено, разделим обе части на $2(1-\rho)$ и выразим длину очереди:
$$
L^{(D)} = \frac{\rho - 2\rho^2 + \E[A_{n+1}^2]}{2(1-\rho)}.
$$

\textbf{3. Вычисление дисперсии приращений.}
Определим $\E[A_{n+1}^2]$ через формулу полной дисперсии:
$$
\mathbb{D}[A_{n+1}] = \E[\mathbb{D}[A_{n+1} \mid S_{n+1}]] + \mathbb{D}[\E[A_{n+1} \mid S_{n+1}]].
$$
Для пуассоновского закона условное математическое ожидание и дисперсия равны $\lambda S_{n+1}$:
$$
\mathbb{D}[A_{n+1}] = \E[\lambda S_{n+1}] + \mathbb{D}[\lambda S_{n+1}] = \lambda \E[S_{n+1}] + \lambda^2 \mathbb{D}[S_{n+1}] = \rho + \lambda^2 \sigma_B^2.
$$
Переходя ко второму начальному моменту:
$$
\E[A_{n+1}^2] = \mathbb{D}[A_{n+1}] + \big(\E[A_{n+1}]\big)^2 = \rho + \lambda^2 \sigma_B^2 + \rho^2.
$$
Подстановка данного выражения в формулу для $L^{(D)}$ дает искомый результат:
$$
L^{(D)} = \frac{\rho - 2\rho^2 + (\rho + \lambda^2 \sigma_B^2 + \rho^2)}{2(1-\rho)} = \frac{2\rho - \rho^2 + \lambda^2 \sigma_B^2}{2(1-\rho)} = \rho + \frac{\rho^2 + \lambda^2 \sigma_B^2}{2(1-\rho)}.
$$

\textbf{4. Совпадение распределений и эквивалентность средних.}

Для завершения доказательства необходимо обосновать равенство $L = L^{(D)}$. Введем три распределения вероятностей состояний системы:
\begin{itemize}
    \item $p_n$ — стационарная вероятность того, что в произвольный момент времени в системе находится $n$ заявок;
    \item $q_n$ — вероятность того, что \textit{прибывающая} заявка застает в системе $n$ заявок;
    \item $\pi_n$ — вероятность того, что \textit{уходящая} после обслуживания заявка оставляет в системе $n$ заявок.
\end{itemize}

Рассмотрим некоторую реализацию процесса $X(t)$ на достаточно большом интервале времени $[0, T]$. В силу того, что приходы и уходы заявок в системе $M|G|1$ являются одиночными событиями, изменения состояния системы происходят только в виде скачков $\pm 1$. 

Зафиксируем уровень $n$ и введем счетчики переходов:
\begin{itemize}
    \item $A_n(T)$ — число переходов вида $n \to n+1$ (приход заявки в систему, находящуюся в состоянии $n$);
    \item $D_n(T)$ — число переходов вида $n+1 \to n$ (уход заявки, оставляющей после себя $n$ заявок).
\end{itemize}

Поскольку для повторного совершения скачка «вверх» система обязана предварительно совершить скачок «вниз» (и наоборот), значения этих счетчиков для любой реализации процесса могут отличаться не более чем на единицу: $|A_n(T) - D_n(T)| \le 1$.

\begin{center}
\begin{tikzpicture}[
    scale=1.3,
    >={Stealth[length=2.5mm]},
    every node/.style={font=\small},
    grid_style/.style={gray!15, thin}
]
    % Сетка
    \draw[grid_style] (0,0) grid (9.5,4);

    % Оси
    \draw[->, thick] (0,0) -- (10,0) node[right] {$t$};
    \draw[->, thick] (0,0) -- (0,4.2) node[above] {$X(t)$};
    
    % Вертикальная ось (состояния)
    \foreach \y/\label in {0.5/n-1, 1.5/n, 2.5/n+1, 3.5/n+2} {
        \draw[thick] (0.05,\y) -- (-0.05,\y) node[left] {$\label$};
    }

    % Горизонтальная ось (моменты событий tau_i)
    \foreach \x in {1,2,...,9} {
        \draw[thick] (\x,0.05) -- (\x,-0.05) node[below] {$\tau_{\x}$};
    }

    % Референтный уровень n
    \draw[thick, DeepSkyBlue4, opacity=0.9] (0,1.5) -- (9.5,1.5) node[right, font=\bfseries\footnotesize] {Уровень $n$};
    \draw[dashed, gray!60] (0,2.5) -- (9.5,2.5);

    % Траектория процесса X(t)
    \draw[ultra thick, Red3] 
        (0,0.5) -- (1,0.5) -- 
        (1,1.5) -- (2,1.5) -- % Приход A_n=1 (D_n=0)
        (2,2.5) -- (3,2.5) -- % Уход D_n=1   (A_n=1)
        (3,1.5) -- (4,1.5) -- % Приход A_n=2 (D_n=1)
        (4,2.5) -- (5,2.5) -- 
        (5,3.5) -- (6,3.5) -- 
        (6,2.5) -- (7,2.5) -- % Уход D_n=2   (A_n=2)
        (7,1.5) -- (8,1.5) -- 
        (8,0.5) -- (9,0.5) -- 
        (9,1.5) -- (9.5,1.5);

    % Обозначения An (Приходы, заставшие n) с текущим Dn
    % tau_2: A_n=1, D_n=0
    \fill[Red3] (2, 1.5) circle (2pt);
    \draw[<-, thick, Red3] (2, 1.6) -- (2, 3.1);
    \node[Red3, fill=white, draw=Red3, thin, rounded corners=1pt, font=\scriptsize\bfseries, align=center] at (2, 3.4) {$A_n = 1$ \\ \color{gray}$D_n = 0$};

    % tau_4: A_n=2, D_n=1
    \fill[Red3] (4, 1.5) circle (2pt);
    \draw[<-, thick, Red3] (4, 1.6) -- (4, 3.1);
    \node[Red3, fill=white, draw=Red3, thin, rounded corners=1pt, font=\scriptsize\bfseries, align=center] at (4, 3.4) {$A_n = 2$ \\ \color{gray}$D_n = 1$};

    % Обозначения Dn (Уходы, оставившие n) с текущим An
    % tau_3: D_n=1, A_n=1
    \fill[DeepSkyBlue4] (3, 2.5) circle (2pt);
    \draw[<-, thick, DeepSkyBlue4] (3, 2.4) -- (3, 0.9);
    \node[DeepSkyBlue4, fill=white, draw=DeepSkyBlue4, thin, rounded corners=1pt, font=\scriptsize\bfseries, align=center] at (3, 0.6) {$D_n = 1$ \\ \color{gray}$A_n = 1$};

    % tau_7: D_n=2, A_n=2
    \fill[DeepSkyBlue4] (7, 2.5) circle (2pt);
    \draw[<-, thick, DeepSkyBlue4] (7, 2.4) -- (7, 0.9);
    \node[DeepSkyBlue4, fill=white, draw=DeepSkyBlue4, thin, rounded corners=1pt, font=\scriptsize\bfseries, align=center] at (7, 0.6) {$D_n = 2$ \\ \color{gray}$A_n = 2$};

\end{tikzpicture}
\end{center}

Пусть $A(T)$ — общее число прибывших заявок, а $D(T)$ — общее число обслуженных заявок за время $T$. Очевидно, что баланс числа заявок в системе задается соотношением:
\begin{equation}\label{eq:balance_total}
    D(T) = A(T) + X(0) - X(T).
\end{equation}

Вычислим вероятность $\pi_n$ как предел отношения частот при $T \to \infty$:
\begin{equation}\label{eq:limit_derivation}
    \pi_n = \lim_{T \to \infty} \frac{D_n(T)}{D(T)} = \lim_{T \to \infty} \frac{A_n(T) + \overbrace{[D_n(T) - A_n(T)]}^{= 0, \pm 1}}{A(T) + \underbrace{X(0) - X(T)}_{\in O(1)}}.
\end{equation}
В стационарном состоянии случайная величина $X(T)$ ограничена (состояние системы не уходит в бесконечность почти наверняка при $\rho < 1$). Следовательно, разность $X(0) - X(T)$ остается ограниченной константой. При $T \to \infty$ общее число поступивших заявок $A(T) \to \infty$. Разделив числитель и знаменатель на $A(T)$, получаем:
\begin{equation*}
    \pi_n = \lim_{T \to \infty} \frac{\frac{A_n(T)}{A(T)} + \frac{0, \pm 1}{A(T)}}{1 + \frac{O(1)}{A(T)}} = \lim_{T \to \infty} \frac{A_n(T)}{A(T)} = q_n.
\end{equation*}

Таким образом, вероятность найти систему в состоянии $n$ в моменты ухода ($\pi_n$) тождественно равна вероятности застать систему в этом же состоянии в моменты прихода ($q_n$). 

Согласно \textbf{свойству PASTA} (Poisson Arrivals See Time Averages), пуассоновский входной поток «видит» систему в стационарном состоянии, что означает $q_n = p_n$. Объединяя результаты, получаем:
\begin{equation*}
    p_n = q_n = \pi_n \quad \implies \quad L = L^{(D)}.
\end{equation*}

Окончательно подставляем выведенное ранее значение $L^{(D)}$ и получаем \textbf{формулу Полячека — Хинчина}:
\begin{equation*}
    \fbox{$\displaystyle L = \rho + \frac{\rho^2 + \lambda^2 \sigma_B^2}{2 (1 - \rho)}$}
\end{equation*}
\end{tproof}

\begin{conseq}
Применение универсального закона Литтла ($L = \lambda W$) к полученному результату позволяет определить все остальные макроскопические метрики системы $M|G|1$:
\begin{itemize}
    \item Среднее время пребывания в системе: $\displaystyle W = \frac{L}{\lambda} = \frac{1}{\mu} + \frac{\rho^2 + \lambda^2 \sigma_B^2}{2\lambda(1-\rho)}$.
    \item Средняя длина очереди: $\displaystyle L_q = L - \rho = \frac{\rho^2 + \lambda^2 \sigma_B^2}{2(1-\rho)}$.
    \item Среднее время ожидания в очереди: $\displaystyle W_q = W - \frac{1}{\mu} = \frac{L_q}{\lambda}$.
\end{itemize}
\end{conseq}

\begin{example}{Сравнение систем M|M|1 и M|D|1}
Формула Полячека — Хинчина позволяет аналитически оценить влияние разброса времени обслуживания на длину очереди. Рассмотрим два случая: систему с показательным временем обслуживания ($M|M|1$) и систему с детерминированным временем обслуживания постоянной длительности ($M|D|1$).
\end{example}

\begin{eproof}

\textbf{Случай 1 (система $M|M|1$).}

Время обслуживания имеет экспоненциальное распределение с параметром $\mu$. Дисперсия экспоненциального распределения равна квадрату математического ожидания: $\sigma_B^2 = \frac{1}{\mu^2}$. Подставляя это в формулу Полячека — Хинчина, выделим слагаемое, отвечающее за очередь:

$$L_{M|M|1} = \rho + \frac{\rho^2 + \lambda^2 \frac{1}{\mu^2}}{2(1-\rho)} = \rho + \frac{2\rho^2}{2(1-\rho)} = \rho + \underbrace{\frac{\rho^2}{1-\rho}}_{L_q}.$$

\textbf{Случай 2 (система $M|D|1$).}

Время обслуживания строго фиксировано: $S = a = \text{const}$. Дисперсия константы равна нулю ($\sigma_B^2 = 0$). При подстановке получаем:

$$L_{M|D|1} = \rho + \frac{\rho^2 + 0}{2(1-\rho)} = \rho + \underbrace{\frac{1}{2} \cdot \frac{\rho^2}{1-\rho}}_{L_q}.$$

Сравнивая вторые слагаемые в обоих случаях, видим, что среднее число заявок в очереди для системы с детерминированным обслуживанием ровно в два раза меньше, чем в простейшей системе с тем же коэффициентом загрузки $\rho$.

Таким образом, устранение стохастичности (вариативности) времени обслуживания позволяет сократить длину очереди вдвое при неизменной интенсивности потока и прибора.
\end{eproof}

\begin{center}
\begin{tikzpicture}
    \begin{axis}[
        width=0.8\textwidth, height=8cm,
        xlabel={Коэффициент загрузки $\rho$}, 
        ylabel={Среднее число заявок $L$},
        xmin=0, xmax=1, ymin=0, ymax=20,
        grid=major,
        legend pos=north west,
        samples=100, thick,
        domain=0:0.99
    ]
        % Формула для M|M|1
        \addplot[Red3, solid] {x/(1-x)};
        \addlegendentry{$M|M|1$ (Экспоненциальное время)}
        
        % Формула для M|D|1
        \addplot[DeepSkyBlue4, dashed] {x + (x^2)/(2*(1-x))};
        \addlegendentry{$M|D|1$ (Детерминированное время)}
        
        % Произвольная система с высокой дисперсией (для наглядности M|G|1)
        % Допустим sigma_B^2 = 3 / mu^2 => lambda^2 sigma_B^2 = 3 rho^2
        \addplot[black, dashdotted] {x + (x^2 + 3*x^2)/(2*(1-x))};
        \addlegendentry{$M|G|1$ (Высокая дисперсия, $\sigma_B^2 = 3/\mu^2$)}
    \end{axis}
\end{tikzpicture}
\captionof*{figure}{\small Зависимость среднего числа заявок в системе от коэффициента загрузки. График демонстрирует, что очереди формируются не только из-за высокой загрузки, но и из-за дисперсии временных параметров.}
\end{center}